\section{결론}

이 글의 결론은 단순하다.
임피던스 제어를 볼 때는 로봇 끝단을 목표점에 묶인 가상 스프링--댐퍼로 먼저 상상하면 된다.
강성은 얼마나 세게 붙잡을지, 감쇠는 움직이는 동안 얼마나 부드럽게 속도를 죽일지, 관성은 외란에 얼마나 둔감하게 반응할지를 정한다.
그래서 좋은 학습 순서는 공식을 한 번에 외우는 것이 아니라, 하나의 파라미터를 바꾸고 응답이 어떻게 달라질지 먼저 예측해 보는 것이다.

이 관점에서 RLC 회로, 질량--스프링--댐퍼 시스템, Panda 가상 벽 실험은 서로 다른 예제가 아니라 같은 질문의 변주이다.
힘과 운동의 관계를 어떻게 설계하고, 그 결과를 시간응답에서 어떻게 읽을 것인가?
MuJoCo 실험실의 \texttt{config.yaml}, \texttt{log.csv}, \texttt{summary.json}, 플롯, \texttt{worksheet.md}는 이 질문을 안전하게 반복해 보는 관찰 도구이다.
여기서 얻는 증거는 실제 하드웨어 접촉 성능의 보증이 아니라, 강성, 감쇠, 목표점, 접촉력 사이의 관계를 손으로 익히기 위한 교육용 관찰값으로 읽으면 된다.

읽고 난 뒤에는 긴 공식 목록보다 몇 가지 감각이 남아야 한다.
원하는 접촉력 \(F\)와 허용 변위 \(\Delta x\)에서 \(k_d \approx F/\Delta x\)로 첫 강성값을 잡고, 허용할 흔들림에 따라 감쇠비를 고른 뒤, \(d_d=2\zeta\sqrt{m_{\mathrm{nom}}k_d}\)가 어떤 의미인지 설명할 수 있으면 된다.
또 위치 오차 하나만으로 접촉 응답을 판단하지 않고, 속도, 계산된 접촉/벽 힘, 제어 입력, effort 지표를 함께 보는 습관이 생기면 이 튜토리얼의 핵심은 이미 잡은 것이다.
조금 더 구체적으로는 다음 기준을 스스로 적용할 수 있으면 좋다.

\begin{itemize}
  \item $1~\mathrm{cm}$의 위치 오차가 강성에 따라 어느 정도의 접촉력으로 바뀌는지 계산할 수 있다.
  \item 원하는 접촉력 $F$와 허용 변위 $\Delta x$에서 $k_d \approx F/\Delta x$라는 첫 강성값을 잡을 수 있다.
  \item 강성을 키우면 정상상태 위치 오차가 줄어드는 대신 접촉력, 포화 가능성, 노이즈 민감도가 커질 수 있음을 설명할 수 있다.
  \item 감쇠비를 바꾸면 오버슈트, 진동, 정착시간이 어떻게 달라지는지 응답 그래프에서 읽을 수 있다.
  \item $d_d=2\zeta\sqrt{m_{\mathrm{nom}}k_d}$처럼 계산한 감쇠계수에서 $m_{\mathrm{nom}}$은 교육용 기준값이며, 실제 끝단 유효질량이나 안전성 보증으로 읽지 않는다.
  \item 접촉 응답을 볼 때 위치 오차 하나만 보지 않고 속도, 이론적 또는 계산된 접촉/벽 힘 신호, 제어 입력, actuator effort 지표, 정의된 가상 에너지, 소산률, 또는 해당 랩에서 제공하는 에너지/effort 지표를 함께 비교할 수 있다.
  \item 로봇 끝단 임피던스, 관절 임피던스, 어드미턴스 제어, 힘 제어가 어떤 상황에서 다른 선택지가 되는지 구분할 수 있다.
  \item 같은 접촉 감각 목표라도 토크/힘 명령이 가능한 로봇에서는 임피던스식 힘/토크 명령으로, 위치 명령만 가능한 로봇에서는 측정/추정 외력 또는 Lab04처럼 계산된 상호작용 신호를 이용해 목표 위치를 바꾸는 어드미턴스식 또는 목표 보정 구현으로 접근함을 구분할 수 있다.
\end{itemize}

이 목록을 모두 동시에 완벽히 할 필요는 없다.
처음에는 세 가지만 확인하면 된다.
목표점에서 얼마나 밀렸는가, 움직이는 동안 얼마나 흔들렸는가, 그리고 그 과정에서 힘이나 effort가 너무 커지지 않았는가.

여기까지 읽었다면 모든 표를 한 번에 외우려고 하지 않아도 된다.
입문자가 이 실험실을 사용할 때의 실제 작업 흐름은 간단하다.
먼저 \texttt{config.yaml}에서 이번에 하나만 바꿀 파라미터를 고르고, 그 변화가 위치, 속도, 힘 또는 effort 플롯을 어떻게 바꿀지 한 문장으로 예측한다.
그다음 실행 후 \texttt{log.csv}, \texttt{summary.json}, \texttt{plots/*.png}를 같은 순서로 열어 설정값, 목표와 입력, 위치와 속도 응답, 힘 또는 제어 입력의 포화 여부를 확인한다.
마지막으로 \texttt{worksheet.md}에 예상과 맞은 점, 달랐던 점, 다음에 하나만 바꿔 볼 값을 적는다.
이 반복이 익숙해지면 임피던스는 어려운 공식 묶음이 아니라, 로봇이 환경을 만났을 때 어떻게 버티고 물러나는지를 읽는 방법이 된다.

실천적으로는 원하는 힘과 허용 변위에서 강성을 잡고, 허용할 흔들림에서 감쇠비를 고른 뒤, 로봇이 받을 수 있는 명령 형식과 사용할 수 있는 힘 신호에 맞춰 임피던스 구현, 어드미턴스식 외부 루프, 또는 목표 보정 구현을 선택한다.
후속 작업에서는 특이점 근처에서 목표 움직임을 무리하지 않게 나누는 Lab03의 감쇠 최소제곱(DLS) 비교와 Lab04의 강성/감쇠 스윕을 더 많은 수업 사례, 학습자 관찰 기록, 정량 지표, 대표 그림으로 확장해 체계화할 필요가 있다.

끝에 남기고 싶은 이미지는 하나다.
로봇 손끝은 목표점에 고무줄과 완충기로 연결되어 있고, 강성은 고무줄의 세기, 감쇠는 흔들림을 죽이는 정도, 관성은 밀렸을 때 바로 움직일지 버틸지를 정한다.
이 그림을 떠올린 뒤 설정을 하나 바꾸고, 그래프에서 그 흔적을 찾는 것이 임피던스 제어를 배우는 가장 좋은 출발점이다.
